\clearpage\thispagestyle{empty}
\pdfbookmark[0]{Quellen zum Brief an Bode}{bode18060730-sources}
\begin{center}\textbf{Quellen zum Brief an Bode}\end{center}
\noindent Diese KI-Übertragung in heutiges Deutsch folgt dem Brief Carl Friedrich Gauß’ vom 30. Juli 1806 an Johann Elert Bode in \href{https://gdz.sub.uni-goettingen.de/id/PPN236007467}{Gauß’ \textit{Werken}, Band~VI} (Göttingen, 1874), S.~278--279. Er erschien in Bodes \textit{Astronomischem Jahrbuch für 1809} (Berlin, 1806), \href{https://www.digitale-sammlungen.de/en/view/bsb10538338?page=225}{S.~215}--\href{https://www.digitale-sammlungen.de/en/view/bsb10538338?page=229}{219}. Die beiden historischen Anmerkungen stehen auch auf \href{https://www.digitale-sammlungen.de/en/view/bsb10538338?page=227}{S.~217 des Jahrbuchs}. Die zweite ist von Bode unterzeichnet.
\par\medskip
Das Jahrbuch druckt die Ephemeride der Pallas auf \href{https://www.digitale-sammlungen.de/en/view/bsb10538338?page=228}{S.~218} und die der Juno auf \href{https://www.digitale-sammlungen.de/en/view/bsb10538338?page=229}{S.~219} ab. Die Tabelle für Juno beginnt dort am 12. Januar 1807. Die Beschreibung in eckigen Klammern in der Werkausgabe nennt den 7. Dezember 1806.
\par\medskip
Die Mailänder Beobachtungen der Juno und der Pallas stehen in der \textit{Monatlichen Correspondenz}, Band~XIV, Juli 1806, auf \href{https://www.digitale-sammlungen.de/en/view/bsb10538606?page=98}{S.~88} bzw.\ \href{https://www.digitale-sammlungen.de/en/view/bsb10538606?page=100}{S.~90}. Die Zeitschrift nennt Carlini als Urheber der Beobachtungen und Berechnungen zur Pallas und der Beobachtungen und Berechnungen zur Juno am Äquatorialsektor. Die historischen Digitalisate stammen von der Bayerischen Staatsbibliothek. Übertragung und Satz wurden mit KI erstellt.
\clearpage
